\section{Agent 时代的工作与就业}

\subsection{数字空间 vs 物理空间}

Karpathy 分析了 BLS（美国劳工统计局）的就业数据，提出了\textbf{数字/物理二元框架}：

\begin{importantbox}{Digital vs Physical}
\begin{itemize}
    \item \textbf{数字空间}：bit 可以复制粘贴，速度接近光速。当前数字世界存在巨大的"unhobbling overhang"——人类没有足够的思考周期来处理已有的数字信息。AI 将在这里引发海量重写和重构
    \item \textbf{物理空间}：原子操作慢得多，资本密集。机器人、自动驾驶等将滞后于数字空间的变革
    \item \textbf{接口层}：数字与物理的界面（传感器/执行器）是下一个大机会——为超级智能提供数据输入和物理世界操作
\end{itemize}
\end{importantbox}

\subsection{Jevons Paradox 与软件需求}

\begin{knowledgebox}{ATM 与银行柜员的经典案例}
ATM 的普及并没有减少银行柜员——反而因为降低了银行网点运营成本，导致更多网点和更多柜员。同样，AI 降低了软件开发成本，可能会\textbf{增加}而非减少软件需求。Karpathy 对软件工程就业持"cautiously optimistic"态度。
\end{knowledgebox}

\subsection{Frontier Labs 的自我替代}

\begin{warningbox}{研究者的悖论}
Frontier Labs 雇佣了约一千名研究者，而这些研究者本质上在"automated themselves away"——主动构建替代自己的自动化系统。Karpathy 回忆在 OpenAI 时曾提醒同事："如果我们成功了，我们所有人都会失业。"
\end{warningbox}

\subsection{为什么软件会成为第一波全面重写的行业}
Karpathy 对软件工程就业的判断并不是简单的乐观或悲观，而是建立在一个更深的前提上：\textbf{软件是目前最适合被 agent 大规模重写的数字系统}。代码天然离散、可复制、可测试、可回滚，很多任务还带有明确的 verifier，这使它成为 auto research 和 code agents 最容易形成飞轮的领域。

\begin{importantbox}{代码是最接近 ``可验证认知劳动'' 的大行业}
\begin{itemize}
    \item 输入、输出和中间状态都可以文本化，适合模型处理；
    \item 许多任务有测试、benchmark、lint、运行日志等客观反馈；
    \item 成本结构以数字资产为主，不像物理行业那样被原子世界拖慢。
\end{itemize}
这解释了为什么 Karpathy 会反复回到 code agents、auto research 和 API 重构，而不是把注意力放在更炫的 humanoid demo 上。
\end{importantbox}

从这个角度看，AI 对软件行业的影响不只是 ``把程序员替换掉一部分''，更像是把原本来不及实现、来不及维护、来不及重构的大量软件需求一次性释放出来。也正因此，Karpathy 才会用 Jevons Paradox 来理解就业问题：效率变高后，最先增长的常常不是裁员速度，而是待实现系统的总量。

\subsection{本章小结}

数字空间将率先经历 AI 带来的剧变，物理空间将滞后跟进。Jevons Paradox 暗示软件需求可能增加。最讽刺的是，AI 研究者自身就是最积极的自我替代者。


